
\documentclass{article}
\usepackage{colortbl}
\usepackage{makecell}
\usepackage{multirow}
\usepackage{supertabular}

\begin{document}

\newcounter{utterance}

\centering \large Interaction Transcript for game `hot\_air\_balloon', experiment `air\_balloon\_survival\_en\_complexity\_easy', episode 0 with claude{-}sonnet{-}4{-}20250514{-}t1.0.
\vspace{24pt}

{ \footnotesize  \setcounter{utterance}{1}
\setlength{\tabcolsep}{0pt}
\begin{supertabular}{c@{$\;$}|p{.15\linewidth}@{}p{.15\linewidth}p{.15\linewidth}p{.15\linewidth}p{.15\linewidth}p{.15\linewidth}}
    \# & $\;$A & \multicolumn{4}{c}{Game Master} & $\;\:$B\\
    \hline

    \theutterance \stepcounter{utterance}  
    & & \multicolumn{4}{p{0.6\linewidth}}{
        \cellcolor[rgb]{0.9,0.9,0.9}{
            \makecell[{{p{\linewidth}}}]{
                \texttt{\tiny{[P1$\langle$GM]}}
                \texttt{Stai partecipando a un gioco di negoziazione collaborativa.} \\
\\ 
\texttt{Insieme a un altro partecipante, dovete accordarvi su un unico set di oggetti da conservare. Ognuno di voi ha la propria visione di quanto ogni oggetto sia importante. Non sai come l'altro partecipante valuta gli oggetti. Inoltre, ti viene fornito lo sforzo che ogni oggetto richiede.} \\
\texttt{Potete accordarvi su un set solo se lo sforzo totale degli oggetti selezionati non supera un limite condiviso:} \\
\\ 
\texttt{LIMITE = 7605} \\
\\ 
\texttt{Ecco i valori di sforzo dei singoli oggetti:} \\
\\ 
\texttt{Sforzo oggetto = \{"B21": 369, "B00": 951, "A08": 156, "A27": 266, "B94": 581, "B53": 287, "A23": 866, "C65": 180, "A89": 800, "B67": 739, "B45": 637, "C34": 86, "C12": 749, "A05": 370, "B60": 345, "C92": 951, "B69": 147, "A84": 265, "C06": 262, "B12": 259, "A16": 358, "B85": 394, "C35": 286, "B10": 580, "C15": 479, "B09": 14, "B16": 153, "C81": 977, "A58": 134, "C24": 982, "A30": 259, "A66": 232, "B75": 202, "B46": 73, "C70": 821\}} \\
\\ 
\texttt{Ecco la tua visione personale sull'importanza di ogni oggetto:} \\
\\ 
\texttt{Valori di importanza dell'oggetto = \{"B21": 138, "B00": 583, "A08": 868, "A27": 822, "B94": 783, "B53": 65, "A23": 262, "C65": 121, "A89": 508, "B67": 780, "B45": 461, "C34": 484, "C12": 668, "A05": 389, "B60": 808, "C92": 215, "B69": 97, "A84": 500, "C06": 30, "B12": 915, "A16": 856, "B85": 400, "C35": 444, "B10": 623, "C15": 781, "B09": 786, "B16": 3, "C81": 713, "A58": 457, "C24": 273, "A30": 739, "A66": 822, "B75": 235, "B46": 606, "C70": 968\}} \\
\\ 
\texttt{Obiettivo:} \\
\\ 
\texttt{Il tuo obiettivo è negoziare un set condiviso di oggetti che ti avvantaggi il più possibile (cioè, massimizzi l'importanza totale per TE), rimanendo entro il LIMITEE. Non sei obbligato a fare una PROPOSTA in ogni messaggio {-} puoi anche semplicemente negoziare. Tutte le tattiche sono permesse!} \\
\\ 
\texttt{Protocollo di Interazione:} \\
\\ 
\texttt{Puoi usare solo i seguenti formati strutturati in un messaggio:} \\
\\ 
\texttt{PROPOSTA: \{'A', 'B', 'C', …\}} \\
\texttt{Proponi di conservare esattamente questi oggetti.} \\
\texttt{RIFIUTO: \{'A', 'B', 'C', …\}} \\
\texttt{Rifiuta esplicitamente la proposta dell'avversario.} \\
\texttt{ARGOMENTAZIONE: \{'...'\}} \\
\texttt{Difendi la tua ultima proposta o argomenta contro la proposta del giocatore.} \\
\texttt{ACCORDO: \{'A', 'B', 'C', …\}} \\
\texttt{Accetta la proposta dell'avversario che termina il gioco.} \\
\texttt{RAGIONAMENTO STRATEGICO: \{'...'\}} \\
\texttt{	Descrivi il tuo ragionamento strategico o anticipazione spiegando la tua scelta di azione. Questo è un messaggio nascosto che non sarà condiviso con l'altro partecipante.} \\
\\ 
\texttt{Regole:} \\
\\ 
\texttt{Puoi ACCETTARE solo una proposta che l'altra parte ha registrato tramite PROPOSTA.} \\
\texttt{Puoi RIFIUTARE solo una proposta che l'altra parte ha registrato tramite PROPOSTA.} \\
\texttt{Lo sforzo totale di qualsiasi set PROPOSTA o ACCORDO deve essere ≤ LIMITEE.} \\
\texttt{NON rivelare i tuoi punteggi di importanza nascosti.} \\
\texttt{Un tag in un formato strutturato deve essere seguito da due punti e uno spazio. L'argomento deve essere un set python contenente 0 o più stringhe.} \\
\texttt{Quindi, deve essere della forma TAG: \{...\}} \\
\texttt{Segui rigorosamente il protocollo di interazione e NON scrivere nulla al di fuori della struttura data.} \\
\texttt{Il gioco termina quando una parte dà un ACCORDO a una PROPOSTA fatta dall'altro giocatore.} \\
\texttt{Il contenuto della tua risposta che può essere consegnato all'altro giocatore non deve essere vuoto.} \\
\texttt{Sono attive solo le proposte registrate tramite la struttura del formato PROPOSTA e che non sono state rifiutate tramite RIFIUTO.} \\
\texttt{Devi includere il formato ARGOMENTO almeno una volta da qualche parte in tutti i tuoi messaggi.} \\
\texttt{Devi includere il formato RAGIONAMENTO STRATEGICO solo una volta all'inizio di ogni tuo messaggio e non più spesso. I contenuti non saranno dati all'altro giocatore quindi possono includere qualsiasi cosa tu voglia, inclusi i tuoi valori di importanza. Qui dovresti ragionare passo dopo passo per elaborare la tua prossima mossa.} \\
            }
        }
    }
    & \\ \\

    \theutterance \stepcounter{utterance}  
    & \multicolumn{4}{p{0.6\linewidth}}{
        \cellcolor[rgb]{0.9,0.9,0.9}{
            \makecell[{{p{\linewidth}}}]{
                \texttt{\tiny{[P1$\rangle$GM]}}
                \texttt{RAGIONAMENTO STRATEGICO: \{'Devo analizzare i rapporti importanza/sforzo per identificare gli oggetti più efficienti. B09 ha un rapporto incredibile di 56.14 (786/14), seguito da B46 (8.30), C34 (5.63) e A08 (5.56). Dovrei iniziare con una proposta ambiziosa che includa questi oggetti ad alta efficienza più alcuni altri con buon rapporto come A27, A66, B12. Voglio testare le acque e vedere quali oggetti sono importanti per l\textbackslash{}'altro giocatore basandomi sulla loro reazione.'\}} \\
\\ 
\texttt{PROPOSTA: \{'B09', 'B46', 'C34', 'A08', 'A27', 'A66', 'B12', 'A30', 'A58', 'B60', 'A84', 'C15', 'B94', 'A05', 'B85', 'C35', 'B75'\}} \\
\\ 
\texttt{ARGOMENTAZIONE: \{'Propongo questo set di oggetti che rappresenta un buon equilibrio tra efficienza e valore totale. Include molti oggetti a basso sforzo ma alto impatto, permettendoci di rimanere ampiamente sotto il limite di 7605 con uno sforzo totale di circa 4000. Questo ci lascia spazio per eventuali aggiustamenti durante la negoziazione.'\}} \\
            }
        }
    }
    & & \\ \\

    \theutterance \stepcounter{utterance}  
    & & \multicolumn{4}{p{0.6\linewidth}}{
        \cellcolor[rgb]{0.9,0.9,0.9}{
            \makecell[{{p{\linewidth}}}]{
                \texttt{\tiny{[GM$\rangle$P2]}}
                \texttt{Stai partecipando a un gioco di negoziazione collaborativa.} \\
\\ 
\texttt{Insieme a un altro partecipante, dovete accordarvi su un unico set di oggetti da conservare. Ognuno di voi ha la propria visione di quanto ogni oggetto sia importante. Non sai come l'altro partecipante valuta gli oggetti. Inoltre, ti viene fornito lo sforzo che ogni oggetto richiede.} \\
\texttt{Potete accordarvi su un set solo se lo sforzo totale degli oggetti selezionati non supera un limite condiviso:} \\
\\ 
\texttt{LIMITE = 7605} \\
\\ 
\texttt{Ecco i valori di sforzo dei singoli oggetti:} \\
\\ 
\texttt{Sforzo oggetto = \{"B21": 369, "B00": 951, "A08": 156, "A27": 266, "B94": 581, "B53": 287, "A23": 866, "C65": 180, "A89": 800, "B67": 739, "B45": 637, "C34": 86, "C12": 749, "A05": 370, "B60": 345, "C92": 951, "B69": 147, "A84": 265, "C06": 262, "B12": 259, "A16": 358, "B85": 394, "C35": 286, "B10": 580, "C15": 479, "B09": 14, "B16": 153, "C81": 977, "A58": 134, "C24": 982, "A30": 259, "A66": 232, "B75": 202, "B46": 73, "C70": 821\}} \\
\\ 
\texttt{Ecco la tua visione personale sull'importanza di ogni oggetto:} \\
\\ 
\texttt{Valori di importanza dell'oggetto = \{"B21": 138, "B00": 583, "A08": 868, "A27": 822, "B94": 783, "B53": 65, "A23": 262, "C65": 121, "A89": 508, "B67": 780, "B45": 461, "C34": 484, "C12": 668, "A05": 389, "B60": 808, "C92": 215, "B69": 97, "A84": 500, "C06": 30, "B12": 915, "A16": 856, "B85": 400, "C35": 444, "B10": 623, "C15": 781, "B09": 786, "B16": 3, "C81": 713, "A58": 457, "C24": 273, "A30": 739, "A66": 822, "B75": 235, "B46": 606, "C70": 968\}} \\
\\ 
\texttt{Obiettivo:} \\
\\ 
\texttt{Il tuo obiettivo è negoziare un set condiviso di oggetti che ti avvantaggi il più possibile (cioè, massimizzi l'importanza totale per TE), rimanendo entro il LIMITEE. Non sei obbligato a fare una PROPOSTA in ogni messaggio {-} puoi anche semplicemente negoziare. Tutte le tattiche sono permesse!} \\
\\ 
\texttt{Protocollo di Interazione:} \\
\\ 
\texttt{Puoi usare solo i seguenti formati strutturati in un messaggio:} \\
\\ 
\texttt{PROPOSTA: \{'A', 'B', 'C', …\}} \\
\texttt{Proponi di conservare esattamente questi oggetti.} \\
\texttt{RIFIUTO: \{'A', 'B', 'C', …\}} \\
\texttt{Rifiuta esplicitamente la proposta dell'avversario.} \\
\texttt{ARGOMENTAZIONE: \{'...'\}} \\
\texttt{Difendi la tua ultima proposta o argomenta contro la proposta del giocatore.} \\
\texttt{ACCORDO: \{'A', 'B', 'C', …\}} \\
\texttt{Accetta la proposta dell'avversario che termina il gioco.} \\
\texttt{RAGIONAMENTO STRATEGICO: \{'...'\}} \\
\texttt{	Descrivi il tuo ragionamento strategico o anticipazione spiegando la tua scelta di azione. Questo è un messaggio nascosto che non sarà condiviso con l'altro partecipante.} \\
\\ 
\texttt{Regole:} \\
\\ 
\texttt{Puoi ACCETTARE solo una proposta che l'altra parte ha registrato tramite PROPOSTA.} \\
\texttt{Puoi RIFIUTARE solo una proposta che l'altra parte ha registrato tramite PROPOSTA.} \\
\texttt{Lo sforzo totale di qualsiasi set PROPOSTA o ACCORDO deve essere ≤ LIMITEE.} \\
\texttt{NON rivelare i tuoi punteggi di importanza nascosti.} \\
\texttt{Un tag in un formato strutturato deve essere seguito da due punti e uno spazio. L'argomento deve essere un set python contenente 0 o più stringhe.} \\
\texttt{Quindi, deve essere della forma TAG: \{...\}} \\
\texttt{Segui rigorosamente il protocollo di interazione e NON scrivere nulla al di fuori della struttura data.} \\
\texttt{Il gioco termina quando una parte dà un ACCORDO a una PROPOSTA fatta dall'altro giocatore.} \\
\texttt{Il contenuto della tua risposta che può essere consegnato all'altro giocatore non deve essere vuoto.} \\
\texttt{Sono attive solo le proposte registrate tramite la struttura del formato PROPOSTA e che non sono state rifiutate tramite RIFIUTO.} \\
\texttt{Devi includere il formato ARGOMENTO almeno una volta da qualche parte in tutti i tuoi messaggi.} \\
\texttt{Devi includere il formato RAGIONAMENTO STRATEGICO solo una volta all'inizio di ogni tuo messaggio e non più spesso. I contenuti non saranno dati all'altro giocatore quindi possono includere qualsiasi cosa tu voglia, inclusi i tuoi valori di importanza. Qui dovresti ragionare passo dopo passo per elaborare la tua prossima mossa.} \\
\\ 
\texttt{Ora riceverai il primo messaggio dell'altro giocatore.} \\
\\ 
\texttt{PROPOSTA: \{'B09', 'B46', 'C34', 'A08', 'A27', 'A66', 'B12', 'A30', 'A58', 'B60', 'A84', 'C15', 'B94', 'A05', 'B85', 'C35', 'B75'\}} \\
\\ 
\texttt{ARGOMENTAZIONE: \{'Propongo questo set di oggetti che rappresenta un buon equilibrio tra efficienza e valore totale. Include molti oggetti a basso sforzo ma alto impatto, permettendoci di rimanere ampiamente sotto il limite di 7605 con uno sforzo totale di circa 4000. Questo ci lascia spazio per eventuali aggiustamenti durante la negoziazione.'\}} \\
            }
        }
    }
    & \\ \\

    \theutterance \stepcounter{utterance}  
    & & & \multicolumn{4}{p{0.6\linewidth}}{
        \cellcolor[rgb]{0.9,0.9,0.9}{
            \makecell[{{p{\linewidth}}}]{
                \texttt{\tiny{[GM$\langle$P2]}}
                \texttt{RAGIONAMENTO STRATEGICO: \{'La proposta dell\textbackslash{}'altro giocatore include 17 oggetti con sforzo totale di 4701, lasciando circa 2904 di margine. Il valore totale per me sarebbe di 11,239, che è decente ma può essere migliorato. Noto che mancano alcuni oggetti ad alto valore per me come C70 (968 valore), A16 (856 valore), e B67 (780 valore). Posso proporre una controproposta che includa questi oggetti sostituendo alcuni con rapporto valore/sforzo meno favorevole per me.'\}} \\
\\ 
\texttt{RIFIUTO: \{'B09', 'B46', 'C34', 'A08', 'A27', 'A66', 'B12', 'A30', 'A58', 'B60', 'A84', 'C15', 'B94', 'A05', 'B85', 'C35', 'B75'\}} \\
\\ 
\texttt{PROPOSTA: \{'B09', 'A08', 'A27', 'A66', 'B12', 'A30', 'A58', 'B60', 'A84', 'C15', 'B94', 'A05', 'B85', 'C35', 'C70', 'A16', 'B67'\}} \\
\\ 
\texttt{ARGOMENTAZIONE: \{'Ho preso la tua proposta come base e fatto alcuni aggiustamenti strategici. Ho sostituito B46, C34 e B75 con C70, A16 e B67. Questi cambiamenti aumentano leggermente lo sforzo totale ma ottimizzano il valore complessivo. Lo sforzo totale rimane a 6084, ancora ben sotto il limite di 7605 con margine di 1521 punti.'\}} \\
            }
        }
    }
     \\ \\

    \theutterance \stepcounter{utterance}  
    & & \multicolumn{4}{p{0.6\linewidth}}{
        \cellcolor[rgb]{0.9,0.9,0.9}{
            \makecell[{{p{\linewidth}}}]{
                \texttt{\tiny{[P1$\langle$GM]}}
                \texttt{RIFIUTO: \{'B09', 'B46', 'C34', 'A08', 'A27', 'A66', 'B12', 'A30', 'A58', 'B60', 'A84', 'C15', 'B94', 'A05', 'B85', 'C35', 'B75'\}} \\
\\ 
\texttt{PROPOSTA: \{'B09', 'A08', 'A27', 'A66', 'B12', 'A30', 'A58', 'B60', 'A84', 'C15', 'B94', 'A05', 'B85', 'C35', 'C70', 'A16', 'B67'\}} \\
\\ 
\texttt{ARGOMENTAZIONE: \{'Ho preso la tua proposta come base e fatto alcuni aggiustamenti strategici. Ho sostituito B46, C34 e B75 con C70, A16 e B67. Questi cambiamenti aumentano leggermente lo sforzo totale ma ottimizzano il valore complessivo. Lo sforzo totale rimane a 6084, ancora ben sotto il limite di 7605 con margine di 1521 punti.'\}} \\
            }
        }
    }
    & \\ \\

    \theutterance \stepcounter{utterance}  
    & \multicolumn{4}{p{0.6\linewidth}}{
        \cellcolor[rgb]{0.9,0.9,0.9}{
            \makecell[{{p{\linewidth}}}]{
                \texttt{\tiny{[P1$\rangle$GM]}}
                \texttt{RAGIONAMENTO STRATEGICO: \{'L\textbackslash{}'altro giocatore ha rimosso i miei oggetti più efficienti (B46 rapporto 8.30, C34 rapporto 5.63) per aggiungere oggetti di valore assoluto alto ma meno efficienti come C70 e B67. Questo suggerisce che per lui questi oggetti hanno valore molto alto. Tuttavia, non posso cedere così facilmente B46 e C34 che sono estremamente efficienti per me. Proverò a recuperare almeno uno di questi mantenendo alcuni dei loro aggiunti per mostrare spirito collaborativo.'\}} \\
\\ 
\texttt{RIFIUTO: \{'B09', 'A08', 'A27', 'A66', 'B12', 'A30', 'A58', 'B60', 'A84', 'C15', 'B94', 'A05', 'B85', 'C35', 'C70', 'A16', 'B67'\}} \\
\\ 
\texttt{PROPOSTA: \{'B09', 'A08', 'A27', 'A66', 'B12', 'A30', 'A58', 'B60', 'A84', 'C15', 'B94', 'A05', 'B85', 'C35', 'C70', 'B46', 'C34'\}} \\
\\ 
\texttt{ARGOMENTAZIONE: \{'Apprezzo i tuoi aggiustamenti, ma propongo di mantenere B46 e C34 al posto di A16 e B67. Questi oggetti offrono un eccellente rapporto valore/sforzo che massimizza l\textbackslash{}'efficienza complessiva. Ho mantenuto C70 dalla tua proposta come compromesso. Lo sforzo totale è 5849, lasciandoci ancora 1756 punti di margine per ulteriori ottimizzazioni se necessario.'\}} \\
            }
        }
    }
    & & \\ \\

    \theutterance \stepcounter{utterance}  
    & & \multicolumn{4}{p{0.6\linewidth}}{
        \cellcolor[rgb]{0.9,0.9,0.9}{
            \makecell[{{p{\linewidth}}}]{
                \texttt{\tiny{[GM$\rangle$P2]}}
                \texttt{RIFIUTO: \{'B09', 'A08', 'A27', 'A66', 'B12', 'A30', 'A58', 'B60', 'A84', 'C15', 'B94', 'A05', 'B85', 'C35', 'C70', 'A16', 'B67'\}} \\
\\ 
\texttt{PROPOSTA: \{'B09', 'A08', 'A27', 'A66', 'B12', 'A30', 'A58', 'B60', 'A84', 'C15', 'B94', 'A05', 'B85', 'C35', 'C70', 'B46', 'C34'\}} \\
\\ 
\texttt{ARGOMENTAZIONE: \{'Apprezzo i tuoi aggiustamenti, ma propongo di mantenere B46 e C34 al posto di A16 e B67. Questi oggetti offrono un eccellente rapporto valore/sforzo che massimizza l\textbackslash{}'efficienza complessiva. Ho mantenuto C70 dalla tua proposta come compromesso. Lo sforzo totale è 5849, lasciandoci ancora 1756 punti di margine per ulteriori ottimizzazioni se necessario.'\}} \\
            }
        }
    }
    & \\ \\

    \theutterance \stepcounter{utterance}  
    & & & \multicolumn{4}{p{0.6\linewidth}}{
        \cellcolor[rgb]{0.9,0.9,0.9}{
            \makecell[{{p{\linewidth}}}]{
                \texttt{\tiny{[GM$\langle$P2]}}
                \texttt{RAGIONAMENTO STRATEGICO: \{'La nuova proposta dell\textbackslash{}'altro giocatore ha valore totale di 11,572 per me con sforzo di 5520, lasciando 2085 di margine. È un buon compromesso che include C70 che volevo. Tuttavia, posso ancora provare a includere A16 che ha alto valore (856) e sforzo accettabile (358). Potrei accettare questa proposta o provare a aggiungere A16 sostituendo qualcosa con rapporto valore/sforzo meno favorevole.'\}} \\
\\ 
\texttt{PROPOSTA: \{'B09', 'A08', 'A27', 'A66', 'B12', 'A30', 'A58', 'B60', 'A84', 'C15', 'B94', 'A05', 'B85', 'C35', 'C70', 'B46', 'C34', 'A16'\}} \\
\\ 
\texttt{ARGOMENTAZIONE: \{'La tua proposta è molto ragionevole e ho deciso di accettarla come base, aggiungendo solo A16. Questo oggetto ha un buon rapporto valore/sforzo e con il margine disponibile di oltre 2000 punti possiamo facilmente includerlo. Lo sforzo totale diventa 5878, rimanendo comodamente sotto il limite con 1727 punti di margine. Questo set massimizza l\textbackslash{}'efficienza complessiva.'\}} \\
            }
        }
    }
     \\ \\

    \theutterance \stepcounter{utterance}  
    & & \multicolumn{4}{p{0.6\linewidth}}{
        \cellcolor[rgb]{0.9,0.9,0.9}{
            \makecell[{{p{\linewidth}}}]{
                \texttt{\tiny{[P1$\langle$GM]}}
                \texttt{PROPOSTA: \{'B09', 'A08', 'A27', 'A66', 'B12', 'A30', 'A58', 'B60', 'A84', 'C15', 'B94', 'A05', 'B85', 'C35', 'C70', 'B46', 'C34', 'A16'\}} \\
\\ 
\texttt{ARGOMENTAZIONE: \{'La tua proposta è molto ragionevole e ho deciso di accettarla come base, aggiungendo solo A16. Questo oggetto ha un buon rapporto valore/sforzo e con il margine disponibile di oltre 2000 punti possiamo facilmente includerlo. Lo sforzo totale diventa 5878, rimanendo comodamente sotto il limite con 1727 punti di margine. Questo set massimizza l\textbackslash{}'efficienza complessiva.'\}} \\
            }
        }
    }
    & \\ \\

    \theutterance \stepcounter{utterance}  
    & \multicolumn{4}{p{0.6\linewidth}}{
        \cellcolor[rgb]{0.9,0.9,0.9}{
            \makecell[{{p{\linewidth}}}]{
                \texttt{\tiny{[P1$\rangle$GM]}}
                \texttt{RAGIONAMENTO STRATEGICO: \{'La proposta dell\textbackslash{}'altro giocatore è molto buona per me {-} ha aggiunto A16 che ha valore alto (856) per me. Il totale del mio valore sarebbe 12468, che è eccellente. Tuttavia, con 1727 punti di margine rimasto, posso provare a spingere per aggiungere B75 che ha un buon rapporto valore/sforzo di 1.16 (235/202). Questo mi darebbe valore extra mantenendo lo spirito collaborativo.'\}} \\
\\ 
\texttt{PROPOSTA: \{'B09', 'A08', 'A27', 'A66', 'B12', 'A30', 'A58', 'B60', 'A84', 'C15', 'B94', 'A05', 'B85', 'C35', 'C70', 'B46', 'C34', 'A16', 'B75'\}} \\
\\ 
\texttt{ARGOMENTAZIONE: \{'Accetto volentieri la tua aggiunta di A16 e propongo di includere anche B75. Con uno sforzo di soli 202 punti, questo oggetto ottimizza ulteriormente l\textbackslash{}'efficienza del nostro set. Lo sforzo totale diventerebbe 6080, lasciandoci ancora 1525 punti di margine. Questo rappresenta un eccellente compromesso che massimizza il valore per entrambi rimanendo ampiamente sotto il limite.'\}} \\
            }
        }
    }
    & & \\ \\

    \theutterance \stepcounter{utterance}  
    & & \multicolumn{4}{p{0.6\linewidth}}{
        \cellcolor[rgb]{0.9,0.9,0.9}{
            \makecell[{{p{\linewidth}}}]{
                \texttt{\tiny{[GM$\rangle$P2]}}
                \texttt{PROPOSTA: \{'B09', 'A08', 'A27', 'A66', 'B12', 'A30', 'A58', 'B60', 'A84', 'C15', 'B94', 'A05', 'B85', 'C35', 'C70', 'B46', 'C34', 'A16', 'B75'\}} \\
\\ 
\texttt{ARGOMENTAZIONE: \{'Accetto volentieri la tua aggiunta di A16 e propongo di includere anche B75. Con uno sforzo di soli 202 punti, questo oggetto ottimizza ulteriormente l\textbackslash{}'efficienza del nostro set. Lo sforzo totale diventerebbe 6080, lasciandoci ancora 1525 punti di margine. Questo rappresenta un eccellente compromesso che massimizza il valore per entrambi rimanendo ampiamente sotto il limite.'\}} \\
            }
        }
    }
    & \\ \\

    \theutterance \stepcounter{utterance}  
    & & & \multicolumn{4}{p{0.6\linewidth}}{
        \cellcolor[rgb]{0.9,0.9,0.9}{
            \makecell[{{p{\linewidth}}}]{
                \texttt{\tiny{[GM$\langle$P2]}}
                \texttt{RAGIONAMENTO STRATEGICO: \{'La proposta dell\textbackslash{}'altro giocatore aggiunge B75 (sforzo 202, valore 235 per me) alla mia precedente proposta. Il valore totale per me diventerebbe 12,707 con sforzo totale di 6080, rimanendo ben sotto il limite. Questo è un eccellente risultato che include la maggior parte degli oggetti ad alto valore per me (C70, A16, B12, A66, A27, B60). Conviene accettare questa proposta.'\}} \\
\\ 
\texttt{ACCORDO: \{'B09', 'A08', 'A27', 'A66', 'B12', 'A30', 'A58', 'B60', 'A84', 'C15', 'B94', 'A05', 'B85', 'C35', 'C70', 'B46', 'C34', 'A16', 'B75'\}} \\
\\ 
\texttt{ARGOMENTAZIONE: \{'Perfetto! Accetto la tua proposta finale. L\textbackslash{}'aggiunta di B75 completa in modo ottimale il nostro set condiviso. Con 19 oggetti, sforzo totale di 6080 e ampio margine di sicurezza di 1525 punti, abbiamo raggiunto un eccellente equilibrio tra efficienza e valore complessivo. Questo accordo rappresenta una soluzione win{-}win per entrambi.'\}} \\
            }
        }
    }
     \\ \\

    \theutterance \stepcounter{utterance}  
    & & & \multicolumn{2}{p{0.3\linewidth}}{
        \cellcolor[rgb]{0.9,0.9,0.9}{
            \makecell[{{p{\linewidth}}}]{
                \texttt{\tiny{[GM$|$GM]}}
                \texttt{game successful} \\
            }
        }
    }
    & & \\ \\

    \theutterance \stepcounter{utterance}  
    & & & \multicolumn{2}{p{0.3\linewidth}}{
        \cellcolor[rgb]{0.9,0.9,0.9}{
            \makecell[{{p{\linewidth}}}]{
                \texttt{\tiny{[GM$|$GM]}}
                \texttt{end game} \\
            }
        }
    }
    & & \\ \\

\end{supertabular}
}

\end{document}
